\makeabstract
    {
    Investigating the Effect of UV Irradiation on Clock Transition-Inducing Defects in Borosilicate Glass
    }
    {
    Ian Moore\textsuperscript{1}, Taylor Hoganson\textsuperscript{1}, Guanchu Chen\textsuperscript{1,2}, Jonathan R. Friedman\textsuperscript{1,2}
    }
    {
    \textsuperscript{1} Department of Physics \& Astronomy, Amherst College, Amherst, MA\\
\textsuperscript{2} University of Massachusetts Amherst, Amherst, MA
    }
    {
    Quantum computing depends on quantum bits, or qubits. Many potential qubits are limited by short coherence times which place an upper bound on the duration of a computation. It is therefore of great importance to find a qubit with a longer coherence time. Coherence time increases dramatically at a clock transition, where the energy difference between two spin states is immune to small fluctuations in magnetic field. Past work has observed a clock transition in borosilicate glass, but the defects causing this transition remain unknown. This project aimed to determine the defects present in borosilicate glass which underlie the clock transition. Previous research by Jeremiah Levy {\textquoteleft}25 indicated that near-UV (300-400 nm wavelength) irradiation could generate clock transition-inducing defects, and this work was intended to refine his findings.
    }
    {
    Using two mercury arc lamps --- one replicating Levy{\textquoteright}s experiment and the other offering greater control --- borosilicate glass samples were irradiated for controlled amounts of time. Defect concentrations within the samples were measured using pulsed electron-spin resonance (ESR) both before and after irradiation and then compared.
    }
    {
    No difference in ESR signal was found after irradiation, and the study was therefore unable to replicate Levy{\textquoteright}s results.
    }
    {
    We thus conclude that near-UV irradiation appears insufficient to generate clock transition-inducing defects in borosilicate glass, suggesting that higher-energy irradiation sources are necessary. This is consistent with the findings of other studies, which show that near-UV irradiation is incapable of generating defects but suggest that far-UV (100-200 nm) irradiation can form significant defects in the glass. While no study currently discusses which defects underlie the clock transition, these studies offer a promising avenue for further research.
    }

    \newpage
    